\documentclass[UTF8,12pt]{ctexart}
\usepackage[a4paper,margin=24mm]{geometry}
\usepackage{amsmath,amssymb,bm,graphicx,adjustbox}
\newcommand{\fitmath}[1]{\begin{adjustbox}{max width=\linewidth}$\displaystyle #1$\end{adjustbox}}
\setlength{\parindent}{0pt}
\setlength{\parskip}{.7em}
\sloppy
\allowdisplaybreaks
\begin{document}
\section*{2011 年考研数学三 · 真题}
\subsection*{原 PDF 第 41 页}

\subsection*{2011年全国硕士研究生招生考试试题}

\subsection*{一、选择题（本题共8小题，每小题4分，共32分。在每小题给出的四个选项中，只有一项符合题目要求，把所选项前的字母填在题后的括号内。）}

（1）已知当 \(\displaystyle x\to0\) 时，函数 \(\displaystyle f(x)=3\sin x-\sin 3x\) 与 \(\displaystyle cx^k\) 是等价无穷小量，则（　　）
（A）\(\displaystyle k=1,\allowbreak c=4\)。 （B）\(\displaystyle k=1,\allowbreak c=-4\)。 （C）\(\displaystyle k=3,\allowbreak c=4\)。 （D）\(\displaystyle k=3,\allowbreak c=-4\)。


（2）设函数 \(\displaystyle f(x)\) 在 \(\displaystyle x=0\) 处可导，且 \(\displaystyle f(0)=0\)，则 \(\displaystyle \lim_{x\to0}\dfrac{\displaystyle x^2f(x)-2f(x^3)}{\displaystyle x^3}=\)（　　）
（A）\(\displaystyle -2f^{\prime}(0)\)。 （B）\(\displaystyle -f^{\prime}(0)\)。 （C）\(\displaystyle f^{\prime}(0)\)。 （D）\(\displaystyle 0\)。


（3）设 \(\displaystyle \{u_n\}\) 是数列，则下列命题正确的是（　　）
（A）若 \(\displaystyle \sum_{n=1}^{\infty}u_n\) 收敛，则 \(\displaystyle \sum_{n=1}^{\infty}(u_{2n-1}+u_{2n})\) 收敛。
（B）若 \(\displaystyle \sum_{n=1}^{\infty}(u_{2n-1}+u_{2n})\) 收敛，则 \(\displaystyle \sum_{n=1}^{\infty}u_n\) 收敛。
（C）若 \(\displaystyle \sum_{n=1}^{\infty}u_n\) 收敛，则 \(\displaystyle \sum_{n=1}^{\infty}(u_{2n-1}-u_{2n})\) 收敛。
（D）若 \(\displaystyle \sum_{n=1}^{\infty}(u_{2n-1}-u_{2n})\) 收敛，则 \(\displaystyle \sum_{n=1}^{\infty}u_n\) 收敛。


（4）设 \(\displaystyle I=\displaystyle \int_0^{\pi/4}\ln(\sin x)\,dx\)，\(\displaystyle J=\displaystyle \int_0^{\pi/4}\ln(\cot x)\,dx\)，\(\displaystyle K=\displaystyle \int_0^{\pi/4}\ln(\cos x)\,dx\)，则 \(\displaystyle I,\allowbreak J,\allowbreak K\) 的大小关系为（　　）
（A）\(\displaystyle I<J<K\)。 （B）\(\displaystyle I<K<J\)。 （C）\(\displaystyle J<I<K\)。 （D）\(\displaystyle K<J<I\)。


（5）设 \(\displaystyle A\) 为3阶矩阵，将 \(\displaystyle A\) 的第2列加到第1列得矩阵 \(\displaystyle B\)，再交换 \(\displaystyle B\) 的第2行与第3行得单位矩阵。记


\[\fitmath{\displaystyle P_1=\begin{pmatrix}1&0&0\\1&1&0\\0&0&1\end{pmatrix},\quad P_2=\begin{pmatrix}1&0&0\\0&0&1\\0&1&0\end{pmatrix}}\]


则 \(\displaystyle A=\)（　　）
（A）\(\displaystyle P_1P_2\)。 （B）\(\displaystyle P_1^{-1}P_2\)。 （C）\(\displaystyle P_2P_1\)。 （D）\(\displaystyle P_2P_1^{-1}\)。


（6）设 \(\displaystyle A\) 为 \(\displaystyle 4\times3\) 矩阵，\(\displaystyle \boldsymbol\eta_1,\allowbreak \boldsymbol\eta_2,\allowbreak \boldsymbol\eta_3\) 是非齐次线性方程组 \(\displaystyle A\boldsymbol x=\boldsymbol\beta\) 的3个线性无关的解，\(\displaystyle k_1,\allowbreak k_2\) 为任意常数，则 \(\displaystyle A\boldsymbol x=\boldsymbol\beta\) 的通解为（　　）
（A）\(\displaystyle \dfrac{\displaystyle \boldsymbol\eta_2+\boldsymbol\eta_3}{\displaystyle 2}+k_1(\boldsymbol\eta_2-\boldsymbol\eta_1)\)。
（B）\(\displaystyle \dfrac{\displaystyle \boldsymbol\eta_2-\boldsymbol\eta_3}{\displaystyle 2}+k_1(\boldsymbol\eta_2-\boldsymbol\eta_1)\)。
（C）\(\displaystyle \dfrac{\displaystyle \boldsymbol\eta_2+\boldsymbol\eta_3}{\displaystyle 2}+k_1(\boldsymbol\eta_2-\boldsymbol\eta_1)+k_2(\boldsymbol\eta_3-\boldsymbol\eta_1)\)。
（D）\(\displaystyle \dfrac{\displaystyle \boldsymbol\eta_2-\boldsymbol\eta_3}{\displaystyle 2}+k_1(\boldsymbol\eta_2-\boldsymbol\eta_1)+k_2(\boldsymbol\eta_3-\boldsymbol\eta_1)\)。


（7）设 \(\displaystyle F_1(x)\) 与 \(\displaystyle F_2(x)\) 为两个分布函数，其相应的概率密度 \(\displaystyle f_1(x)\) 与 \(\displaystyle f_2(x)\) 是连续函数，则必为概率密度的是（　　）
（A）\(\displaystyle f_1(x)f_2(x)\)。 （B）\(\displaystyle 2f_2(x)F_1(x)\)。 （C）\(\displaystyle f_1(x)F_2(x)\)。 （D）\(\displaystyle f_1(x)F_2(x)+f_2(x)F_1(x)\)。


\clearpage

\subsection*{原 PDF 第 42 页}

（8）设总体 \(\displaystyle X\) 服从参数为 \(\displaystyle \lambda\;(\lambda>0)\) 的泊松分布，\(\displaystyle X_1,\allowbreak X_2,\allowbreak \ldots,\allowbreak X_n\;(n\ge2)\) 为来自该总体的简单随机样本，则对于统计量 \(\displaystyle T_1=\dfrac{\displaystyle 1}{\displaystyle n}\displaystyle \sum_{i=1}^{n}X_i\) 和 \(\displaystyle T_2=\dfrac{\displaystyle 1}{\displaystyle n-1}\displaystyle \sum_{i=1}^{n-1}X_i+\dfrac{\displaystyle 1}{\displaystyle n}X_n\)，有（　　）
（A）\(\displaystyle E(T_1)>E(T_2),\allowbreak D(T_1)>D(T_2)\)。 （B）\(\displaystyle E(T_1)>E(T_2),\allowbreak D(T_1)<D(T_2)\)。
（C）\(\displaystyle E(T_1)<E(T_2),\allowbreak D(T_1)>D(T_2)\)。 （D）\(\displaystyle E(T_1)<E(T_2),\allowbreak D(T_1)<D(T_2)\)。


\subsection*{二、填空题（本题共6小题，每小题4分，共24分，把答案填在题中横线上。）}

（9）设 \(\displaystyle f(x)=\displaystyle \lim_{t\to0}x(1+3t)^{x/t}\)，则 \(\displaystyle f^{\prime}(x)=\)\_\_\_\_\_\_。


（10）设函数 \(\displaystyle z=\left(1+\dfrac{\displaystyle x}{\displaystyle y}\right)^{x/y}\)，则 \(\displaystyle dz\big|_{(1,\allowbreak 1)}=\)\_\_\_\_\_\_。


（11）曲线 \(\displaystyle \tan\left(x+y+\dfrac{\displaystyle \pi}{\displaystyle 4}\right)=e^y\) 在点 \(\displaystyle (0,\allowbreak 0)\) 处的切线方程为\_\_\_\_\_\_。


（12）曲线 \(\displaystyle y=\sqrt{x^2-1}\)，直线 \(\displaystyle x=2\) 及 \(\displaystyle x\) 轴所围的平面图形绕 \(\displaystyle x\) 轴旋转所成的旋转体的体积为\_\_\_\_\_\_。


（13）设二次型 \(\displaystyle f(x_1,\allowbreak x_2,\allowbreak x_3)=\boldsymbol x^{\mathrm T}A\boldsymbol x\) 的秩为1，\(\displaystyle A\) 的各行元素之和为3，则 \(\displaystyle f\) 在正交变换 \(\displaystyle \boldsymbol x=Q\boldsymbol y\) 下的标准形为\_\_\_\_\_\_。


（14）设二维随机变量 \(\displaystyle (X,\allowbreak Y)\) 服从正态分布 \(\displaystyle N(\mu,\allowbreak \mu;\sigma^2,\allowbreak \sigma^2;0)\)，则 \(\displaystyle E(XY^2)=\)\_\_\_\_\_\_。


\subsection*{三、解答题（本题共9小题，共94分，解答应写出文字说明、证明过程或演算步骤。）}

（15）（本题满分10分）求极限 \(\displaystyle \lim_{x\to0}\dfrac{\displaystyle \sqrt{1+2\sin x}-x-1}{\displaystyle x\ln(1+x)}\)。


（16）（本题满分10分）已知函数 \(\displaystyle f(u,\allowbreak v)\) 具有二阶连续偏导数，\(\displaystyle f(1,\allowbreak 1)=2\) 是 \(\displaystyle f(u,\allowbreak v)\) 的极值，\(\displaystyle z=f(x+y,\allowbreak f(x,\allowbreak y))\)。求 \(\displaystyle \left.\dfrac{\displaystyle \partial^2 z}{\displaystyle \partial x\partial y}\right|_{(1,\allowbreak 1)}\)。


\clearpage

\subsection*{原 PDF 第 43 页}

（17）（本题满分10分）求不定积分 \(\displaystyle \int\dfrac{\displaystyle \arcsin\sqrt{x}+\ln x}{\displaystyle \sqrt{x}}\,dx\)。


（18）（本题满分10分）证明方程 \(\displaystyle 4\arctan x-x+\dfrac{\displaystyle 4\pi}{\displaystyle 3}-\sqrt3=0\) 恰有两个实根。


（19）（本题满分10分）设函数 \(\displaystyle f(x)\) 在区间 \(\displaystyle [0,\allowbreak 1]\) 上具有连续导数，\(\displaystyle f(0)=1\)，且满足


\[\fitmath{\displaystyle \iint_{D_t}f^{\prime}(x+y)\,dx\,dy=\displaystyle \iint_{D_t}f(t)\,dx\,dy,\quad D_t=\{(x,y)\mid 0\le y\le t-x,\ 0\le x\le t\}\ (0<t\le1).}\]


求 \(\displaystyle f(x)\) 的表达式。


（20）（本题满分11分）设向量组 \(\displaystyle \boldsymbol\alpha_1=(1,\allowbreak 0,\allowbreak 1)^{\mathrm T},\allowbreak \boldsymbol\alpha_2=(0,\allowbreak 1,\allowbreak 1)^{\mathrm T},\allowbreak \boldsymbol\alpha_3=(1,\allowbreak 3,\allowbreak 5)^{\mathrm T}\) 不能由向量组 \(\displaystyle \boldsymbol\beta_1=(1,\allowbreak 1,\allowbreak 1)^{\mathrm T},\allowbreak \boldsymbol\beta_2=(1,\allowbreak 2,\allowbreak 3)^{\mathrm T},\allowbreak \boldsymbol\beta_3=(3,\allowbreak 4,\allowbreak a)^{\mathrm T}\) 线性表示。
（I）求 \(\displaystyle a\) 的值；
（II）将 \(\displaystyle \boldsymbol\beta_1,\allowbreak \boldsymbol\beta_2,\allowbreak \boldsymbol\beta_3\) 用 \(\displaystyle \boldsymbol\alpha_1,\allowbreak \boldsymbol\alpha_2,\allowbreak \boldsymbol\alpha_3\) 线性表示。


\clearpage

\subsection*{原 PDF 第 44 页}

（21）（本题满分11分）设 \(\displaystyle A\) 为3阶实对称矩阵，\(\displaystyle A\) 的秩为2，且


\[\fitmath{\displaystyle A\begin{pmatrix}1&1\\0&0\\-1&1\end{pmatrix}=\begin{pmatrix}-1&1\\0&0\\1&1\end{pmatrix}.}\]


（I）求 \(\displaystyle A\) 的所有特征值与特征向量；
（II）求矩阵 \(\displaystyle A\)。


（22）（本题满分11分）设随机变量 \(\displaystyle X\) 与 \(\displaystyle Y\) 的概率分布分别为


\[\fitmath{\displaystyle \begin{array}{c|cc}X&0&1\\\hline P&\dfrac{\displaystyle 1}{\displaystyle 3}&\dfrac{\displaystyle 2}{\displaystyle 3}\end{array}\qquad\begin{array}{c|ccc}Y&-1&0&1\\\hline P&\dfrac{\displaystyle 1}{\displaystyle 3}&\dfrac{\displaystyle 1}{\displaystyle 3}&\dfrac{\displaystyle 1}{\displaystyle 3}\end{array}}\]


且 \(\displaystyle P\{X^2=Y^2\}=1\)。
（I）求二维随机变量 \(\displaystyle (X,\allowbreak Y)\) 的概率分布；
（II）求 \(\displaystyle Z=XY\) 的概率分布；
（III）求 \(\displaystyle X\) 与 \(\displaystyle Y\) 的相关系数 \(\displaystyle \rho_{XY}\)。


（23）（本题满分11分）设二维随机变量 \(\displaystyle (X,\allowbreak Y)\) 服从区域 \(\displaystyle G\) 上的均匀分布，其中 \(\displaystyle G\) 是由 \(\displaystyle x-y=0\)，\(\displaystyle x+y=2\) 与 \(\displaystyle y=0\) 所围成的三角形区域。
（I）求 \(\displaystyle X\) 的概率密度 \(\displaystyle f_X(x)\)；
（II）求条件概率密度 \(\displaystyle f_{X\mid Y}(x\mid y)\)。


\clearpage
\end{document}
